\documentclass{beamer}
\usetheme{Madrid}
\usecolortheme{default}
\usepackage{listings}
\usepackage{color}
\usepackage{xcolor}
\usepackage{amsmath}
\usepackage{amssymb}
\usepackage{graphicx}
\usepackage{booktabs}

% Define code style
\lstset{
    language=Python,
    basicstyle=\ttfamily\small,
    keywordstyle=\color{blue},
    commentstyle=\color{gray},
    stringstyle=\color{red},
    breaklines=true,
    showstringspaces=false,
    tabsize=2,
    backgroundcolor=\color{lightgray!20}
}

\title{Week 6: From Random Forest to LightGBM}
\subtitle{Mastering Gradient Boosting for Medical Data}
\author{MDST Project: Heart Failure Survival Analysis}
\date{2026}

\begin{document}

%=============================================================
% TITLE SLIDE
%=============================================================
\frame{\titlepage}

%=============================================================
% SLIDE 1: Learning Objectives
%=============================================================
\begin{frame}[fragile]
\frametitle{Learning Objectives}

By the end of this week, you will:

\begin{itemize}
    \item[\checkmark] Understand \textbf{Random Forest baseline} \\
    \quad (parallel independent trees from Week 5)

    \vspace{0.3cm}
    \item[\checkmark] Learn \textbf{Gradient Boosting concept} \\
    \quad (sequential error correction)

    \vspace{0.3cm}
    \item[\checkmark] Master \textbf{LightGBM} \\
    \quad (modern leaf-wise boosting with hyperparameter optimization)

    \vspace{0.3cm}
    \item[\checkmark] Evaluate using \textbf{4 metrics} \\
    \quad (Accuracy, F1-Score, ROC-AUC, MCC)

    \vspace{0.3cm}
    \item[\checkmark] Get \textbf{hands-on practice} \\
    \quad (hyperparameter tuning with Optuna)
\end{itemize}

\end{frame}

%=============================================================
% SLIDE 2: The Four Metrics
%=============================================================
\begin{frame}[fragile]
\frametitle{Evaluation Metrics: Following Chicco \& Jurman (2020)}

\vspace{-0.3cm}
{\small
\begin{tabular}{|p{1.5cm}|p{2.5cm}|p{2.5cm}|p{1.8cm}|}
\hline
\textbf{Metric} & \textbf{Formula} & \textbf{Interpretation} & \textbf{Best For} \\
\hline
Accuracy & $\frac{TP+TN}{TP+TN+FP+FN}$ & \% correct & Balanced \\
\hline
F1-Score & $2 \cdot \frac{P \cdot R}{P+R}$ & Precision-Recall & Imbalanced \\
\hline
ROC-AUC & Area under curve & Ranking quality & Threshold-free \\
\hline
MCC & Correlation coeff. & Most robust & Binary imbalanced \\
\hline
\end{tabular}
}

\vspace{0.5cm}
\textbf{Key Insight:} No single metric tells the full story.

\vspace{0.3cm}
\begin{center}
    \textit{Report all four metrics for comprehensive evaluation}
\end{center}

\end{frame}

%=============================================================
% SLIDE 3: Why Multiple Metrics Matter
%=============================================================
\begin{frame}[fragile]
\frametitle{Why Multiple Metrics Matter: A Real Example}

\textbf{Scenario:} Heart failure classification (10\% positive class)

\vspace{0.3cm}
\begin{center}
{\small
\begin{tabular}{|l|c|c|c|c|}
\hline
\textbf{Model} & \textbf{Accuracy} & \textbf{F1} & \textbf{ROC-AUC} & \textbf{MCC} \\
\hline
A: Always predict ``No'' & 90\% & 0\% & 50\% & 0\% \\
\hline
B: Balanced predictions & 75\% & 60\% & 85\% & 50\% \\
\hline
\end{tabular}
}
\end{center}

\vspace{0.4cm}
\textbf{The Problem:} Model A looks excellent by accuracy alone!

\vspace{0.2cm}
\textbf{The Solution:} F1, ROC-AUC, and MCC expose the failure.

\vspace{0.3cm}
\begin{alertblock}{Lesson}
Accuracy alone is \textbf{misleading} on imbalanced data. \\
Use MCC and ROC-AUC as primary metrics.
\end{alertblock}

\end{frame}

%=============================================================
% SLIDE 4: Random Forest Recap
%=============================================================
\begin{frame}[fragile]
\frametitle{Random Forest: Parallel Learning (Week 5 Review)}

\textbf{How It Works:}

\begin{center}
\begin{tabular}{c}
Bootstrap sampling (random rows with replacement) \\
$\downarrow$ \\
Train trees independently on different samples \\
$\downarrow$ \\
Each tree sees random subset of features \\
$\downarrow$ \\
Majority vote on final prediction \\
\end{tabular}
\end{center}

\vspace{0.4cm}
\textbf{Key Characteristics:}

\begin{itemize}
    \item \textbf{Process:} Parallel (all trees trained simultaneously)
    \item \textbf{Tree depth:} Deep (10-12 levels)
    \item \textbf{Learning mechanism:} Variance reduction via averaging
    \item \textbf{Speed:} Fast training
    \item \textbf{Role:} Good baseline for Week 5
\end{itemize}

\end{frame}

%=============================================================
% SLIDE 5: RF In Action
%=============================================================
\begin{frame}[fragile]
\frametitle{Random Forest In Action: Voting Example}

\textbf{Example with 4 trees:}

\vspace{0.2cm}
{\scriptsize
\begin{lstlisting}
Original data: [row1, row2, row3, row4, row5]

Tree 1 (bootstrap sample): [row1, row1, row3, row5, row5] → Yes
Tree 2 (bootstrap sample): [row2, row3, row3, row4, row5] → No
Tree 3 (bootstrap sample): [row1, row2, row4, row4, row5] → Yes
Tree 4 (bootstrap sample): [row2, row3, row4, row5, row5] → Yes

Final Prediction: Yes  (3/4 trees agree = majority vote)
\end{lstlisting}
}

\vspace{0.3cm}
\textbf{Why Averaging Works:}

\begin{itemize}
    \item Each tree is a \textbf{high-variance} estimate (overfits to its sample)
    \item Errors are \textbf{uncorrelated} across samples
    \item Ensemble average has \textbf{lower variance} than individual trees
    \item Mathematical insight: $\text{Var}(\text{average}) = \frac{\text{Var}}{n}$
\end{itemize}

\end{frame}

%=============================================================
% SLIDE 6: Gradient Boosting Paradigm
%=============================================================
\begin{frame}[fragile]
\frametitle{Gradient Boosting: Sequential Learning Paradigm}

\textbf{The Sequential Idea:}

\begin{center}
{\small
\begin{tabular}{l}
Iteration 1: Train Tree$_1$ on original data \\
\\
Iteration 2: Compute residuals (errors) from Tree$_1$ \\
\quad\quad\quad\quad Train Tree$_2$ on these residuals \\
\\
Iteration 3: Compute residuals from Tree$_1$ + Tree$_2$ \\
\quad\quad\quad\quad Train Tree$_3$ on these residuals \\
\\
\ldots repeat \ldots \\
\\
Final Prediction: Tree$_1$ + $\alpha$ Tree$_2$ + $\alpha$ Tree$_3$ + \ldots
\end{tabular}
}
\end{center}

\vspace{0.3cm}
\begin{center}
\textbf{\textcolor{red}{Each new tree focuses on fixing mistakes from previous trees}}
\end{center}

\end{frame}

%=============================================================
% SLIDE 7: Why GB Works Better
%=============================================================
\begin{frame}[fragile]
\frametitle{Why Gradient Boosting Outperforms Random Forest}

\begin{center}
{\small
\begin{tabular}{|l|c|c|}
\hline
\textbf{Aspect} & \textbf{Random Forest} & \textbf{Gradient Boosting} \\
\hline
\textbf{Bias Reduction} & No & \textcolor{green}{Yes} \\
\hline
\textbf{Variance Reduction} & \textcolor{green}{Yes} & \textcolor{green}{Yes} \\
\hline
\textbf{Tree Depth} & Deep (10-12) & Shallow (3-5) \\
\hline
\textbf{Can Learn Errors} & No & \textcolor{green}{Yes} \\
\hline
\textbf{Typical Accuracy} & $\sim$80-85\% & $\sim$85-90\% \\
\hline
\end{tabular}
}
\end{center}

\vspace{0.4cm}
\textbf{Mathematical Insight:}

\begin{itemize}
    \item \textbf{Random Forest:} $\text{Error} = \text{Bias}^2 + \text{Variance}$
    \begin{itemize}
        \item Reduces variance only
        \item Asymptotically: Error $\approx$ Bias$^2$ (bottleneck)
    \end{itemize}

    \item \textbf{Gradient Boosting:} $\text{Error} = \text{Bias}^2 + \text{Variance}$
    \begin{itemize}
        \item Reduces \textbf{both} bias and variance
        \item Better asymptotic error
    \end{itemize}
\end{itemize}

\end{frame}

%=============================================================
% SLIDE 8: GB Learning From Residuals
%=============================================================
\begin{frame}[fragile]
\frametitle{Gradient Boosting Example: Learning from Residuals}

\textbf{Realistic Scenario:}

\vspace{0.2cm}
{\small
\begin{lstlisting}
Data samples: [(age=50, died=No),
               (age=60, died=Yes),
               (age=55, died=No)]

Tree 1 predictions:  [0.0,  0.8,  0.2]
Actual labels:       [0.0,  1.0,  0.0]
Residuals:          [0.0, +0.2, -0.2]  ← Tree 2 focuses here

Tree 2 learns:
  - For sample 2: Predict +0.2 (boost the prediction)
  - For sample 3: Predict -0.2 (reduce the prediction)

After Tree 1 + 0.1 × Tree 2:
  Predictions: [0.0, 0.82, 0.18]
  ↓
  Closer to actual [0.0, 1.0, 0.0]

(Repeat for Tree 3, Tree 4, ...)
\end{lstlisting}
}

\vspace{0.2cm}
\textbf{Key Insight:} Each tree focuses on \textit{correcting} previous mistakes.

\end{frame}

%=============================================================
% SLIDE 9: GB Hyperparameters
%=============================================================
\begin{frame}[fragile]
\frametitle{Gradient Boosting Hyperparameters Explained}

\textbf{n\_estimators} (number of trees)
\begin{itemize}
    \item More trees = more learning iterations
    \item Trade-off: More complexity, slower training
    \item Typical range: 50 to 200
\end{itemize}

\vspace{0.2cm}
\textbf{learning\_rate} (step size, also called `shrinkage')
\begin{itemize}
    \item lr = 0.1: Each tree contributes 10\% to prediction
    \item lr = 0.05: Each tree contributes 5\% (slower, more stable)
    \item lr = 0.01: Each tree contributes 1\% (very slow, careful learning)
    \item Rule of thumb: Lower lr needs more trees
\end{itemize}

\vspace{0.2cm}
\textbf{max\_depth} (tree complexity)
\begin{itemize}
    \item depth = 3: Shallow trees (avoid overfitting)
    \item depth = 5: Medium trees (balance)
    \item depth = 10: Deep trees (high variance, prone to overfitting)
\end{itemize}

\vspace{0.2cm}
\textbf{subsample} (data fraction per tree)
\begin{itemize}
    \item 0.8 = use 80\% of data for each tree
    \item Adds randomness, prevents overfitting
\end{itemize}

\end{frame}

%=============================================================
% SLIDE 10: GB Training Intuition
%=============================================================
\begin{frame}[fragile]
\frametitle{Gradient Boosting Training Curve Intuition}

\textbf{Learning trajectory on training data:}

\vspace{0.2cm}
{\small
\begin{lstlisting}
Iteration 1: Train on full data
  Misclassification rate: 25%

Iteration 2: Train on residuals of 25% samples
  New misclassification: 20%

Iteration 3: Train on residuals of remaining 20%
  New misclassification: 18%

Iteration 4: Train on residuals of remaining 18%
  New misclassification: 17%

Iteration 5: Train on residuals of remaining 17%
  New misclassification: 16.5%

(Diminishing returns as we improve)
\end{lstlisting}
}

\vspace{0.2cm}
\textbf{Pattern:} Learning curve is \textbf{steep at first}, then \textbf{flattens}.

\vspace{0.2cm}
This is why we typically use 50-100 trees (good balance between improvement and computation).

\end{frame}

%=============================================================
% SLIDE 11: LightGBM Introduction
%=============================================================
\begin{frame}[fragile]
\frametitle{LightGBM: The Modern Upgrade}

\textbf{Traditional Gradient Boosting (Level-wise):}

All nodes at depth $d$ split before any node at depth $d+1$.

\vspace{0.2cm}
\begin{center}
{\tiny
\begin{verbatim}
           node
          /    \
        node    node      <- depth 1
       /  \    /   \
     node node node node   <- depth 2
     /\   /\   /\   /\
    l  l l  l l  l l  l    <- depth 3 (ALL split)
\end{verbatim}
}
\end{center}

\vspace{0.2cm}
\textbf{LightGBM (Leaf-wise):}

Splits only the leaf that reduces loss most (smart focus).

\vspace{0.2cm}
\begin{center}
{\tiny
\begin{verbatim}
           node
          /    \
        node    node
       /  \      |
     node node  node   <- splits only where it matters most
\end{verbatim}
}
\end{center}

\vspace{0.2cm}
\textbf{Result:} Fewer trees needed, \textbf{10-20× faster}, often better accuracy.

\end{frame}

%=============================================================
% SLIDE 12: LightGBM vs Competitors
%=============================================================
\begin{frame}[fragile]
\frametitle{LightGBM vs Other Boosting Algorithms}

\begin{center}
{\small
\begin{tabular}{|l|c|c|c|c|}
\hline
\textbf{Algorithm} & \textbf{Growth} & \textbf{Speed} & \textbf{Accuracy} & \textbf{Overfit} \\
\hline
GB (sklearn) & Level-wise & Baseline & Good & Low \\
\hline
XGBoost & Level-wise & Medium & Very Good & Medium \\
\hline
\textbf{LightGBM} & \textbf{Leaf-wise} & \textbf{Very Fast} & \textbf{Excellent} & \textbf{High*} \\
\hline
CatBoost & Symmetric & Slow & Excellent & Medium \\
\hline
\end{tabular}
}
\end{center}

\vspace{0.3cm}
\textit{* High overfitting risk on small data; requires careful regularization}

\vspace{0.3cm}
\textbf{For our dataset:} Heart failure data has only 299 samples.

\begin{itemize}
    \item LightGBM with strong regularization is optimal
    \item Conservative hyperparameters prevent overfitting
    \item Benefit: 10× speedup vs sklearn GB
\end{itemize}

\end{frame}

%=============================================================
% SLIDE 13: LightGBM Hyperparameters
%=============================================================
\begin{frame}[fragile]
\frametitle{LightGBM Hyperparameters}

\textbf{num\_leaves} (tree complexity limit)
\begin{itemize}
    \item 15: Conservative (prevents overfitting)
    \item 31: Default, balanced
    \item 100: Aggressive (high variance, risk on small data)
\end{itemize}

\textbf{learning\_rate}
\begin{itemize}
    \item Same interpretation as GB
    \item Default: 0.1 (standard), 0.05 (slower, stable)
\end{itemize}

\textbf{min\_data\_in\_leaf}
\begin{itemize}
    \item Minimum samples required in each leaf
    \item Default: 20
    \item Heart failure data: Use 20-30 as safety valve against overfitting
\end{itemize}

\textbf{feature\_fraction \& bagging\_fraction}
\begin{itemize}
    \item Randomness in features and samples
    \item Default: 1.0 (use all)
    \item Recommended: 0.8 (use 80\%)
\end{itemize}

\textbf{max\_depth}
\begin{itemize}
    \item Explicitly limit depth (safety mechanism)
    \item Useful: max\_depth=7 (overrides num\_leaves if needed)
\end{itemize}

\end{frame}

%=============================================================
% SLIDE 14: Leaf-wise Growth Advantage
%=============================================================
\begin{frame}[fragile]
\frametitle{Why Leaf-wise Growth is Smarter}

\textbf{The Smart Decision:}

\vspace{0.2cm}
{\small
\begin{lstlisting}
Training data analysis:
  Node A: 100 samples with 30% error rate
  Node B: 80 samples with 5% error rate

Level-wise (traditional):
  Split both A and B equally
  → Wastes effort on node B

Leaf-wise (LightGBM):
  Split A first (30% error is higher!)
  → Fix the bigger problem first
  → Then split B (5% error is now bigger)

Result: Finds high-error regions faster
        with fewer total splits needed
\end{lstlisting}
}

\vspace{0.3cm}
\textbf{The Danger:}

On small data (n=299), leaf-wise can find \textbf{spurious patterns}.

\vspace{0.2cm}
\begin{alertblock}{Safety Mechanism}
Use \texttt{min\_data\_in\_leaf=20} as a safety valve to prevent overfitting.
\end{alertblock}

\end{frame}

%=============================================================
% SLIDE 15: Optuna Motivation
%=============================================================
\begin{frame}[fragile]
\frametitle{Why Optuna for Hyperparameter Tuning?}

\textbf{Random Search:}
\begin{itemize}
    \item Try 50 random combinations
    \item Expected good params: Maybe 2-3 by luck
    \item Wasteful
\end{itemize}

\vspace{0.2cm}
\textbf{Grid Search:}
\begin{itemize}
    \item Try $3 \times 3 \times 3 \times 3 = 81$ combinations
    \item Computational cost: Massive
    \item Covers only discrete points
\end{itemize}

\vspace{0.2cm}
\textbf{Bayesian Optimization (Optuna):}
\begin{itemize}
    \item Try 50 trials, but \textbf{smart}
    \item Trials 1-10: Explore broadly
    \item Trials 11-30: Focus on promising regions
    \item Trials 31-50: Refine around the best
    \item Expected good params: 10-20
\end{itemize}

\vspace{0.3cm}
\begin{center}
\textbf{Optuna wins:} Best quality with fixed computational budget.
\end{center}

\end{frame}

%=============================================================
% SLIDE 16: Optuna Strategy
%=============================================================
\begin{frame}[fragile]
\frametitle{Optuna Strategy for Our Model}

\textbf{The optimization function:}

{\small
\begin{lstlisting}
def objective(trial):
    # Suggest hyperparameters
    n_est = trial.suggest_int('n_estimators', 50, 150)
    lr = trial.suggest_float('learning_rate', 0.01, 0.1, log=True)
    num_leaves = trial.suggest_int('num_leaves', 15, 50)
    # ... more hyperparameters ...

    # Train and evaluate
    model = LGBMClassifier(**params)
    scores = cross_val_score(model, X, y, cv=5,
                            scoring='roc_auc')

    # Return single metric for optimization
    return scores.mean()

# Let Optuna run 50 trials
study.optimize(objective, n_trials=50)
\end{lstlisting}
}

\vspace{0.2cm}
\textbf{Key Decisions in Our Setup:}

\begin{itemize}
    \item \textbf{log-scale for learning\_rate:} Makes 0.001, 0.01, 0.1, 0.3 equally spaced
    \item \textbf{ROC-AUC metric:} Handles imbalanced data better
    \item \textbf{5-fold CV:} Honest, robust performance estimate
\end{itemize}

\end{frame}

%=============================================================
% SLIDE 17: Full Pipeline
%=============================================================
\begin{frame}[fragile]
\frametitle{The Complete ML Pipeline}

\vspace{-0.3cm}
\begin{center}
{\scriptsize
\begin{tabular}{c}
Raw Data (299 samples, 11 features) \\
$\downarrow$ \\
Train/Test Split (70/30) + StandardScaler \\
$\downarrow$ \\
\end{tabular}

\vspace{0.1cm}
\begin{tabular}{cc}
Training Set (210) & $\rightarrow$ & Test Set (89) \\
(use only for training) & & (final evaluation)
\end{tabular}

\vspace{0.1cm}
\begin{tabular}{c}
Training Set \\
$\downarrow$ \\
\begin{tabular}{|c|}
\hline
\text{Section 2: Random Forest (baseline)} \\
\hline
\text{Section 4: Gradient Boosting} \\
\hline
\text{Section 5: LightGBM (simple)} \\
\hline
\text{Section 6: Optuna Tuning (50 trials)} \\
\hline
\text{Section 6: LightGBM (tuned)} \\
\hline
\end{tabular} \\
$\downarrow$ \\
All models evaluated on Test Set \\
$\downarrow$ \\
Compare 4 metrics: Accuracy, F1, ROC-AUC, MCC
\end{tabular}
}
\end{center}

\end{frame}

%=============================================================
% SLIDE 18: Expected Results
%=============================================================
\begin{frame}[fragile]
\frametitle{Expected Results on Heart Failure Data}

\textbf{Typical Performance:}

\begin{center}
{\small
\begin{tabular}{|l|c|c|c|c|}
\hline
\textbf{Model} & \textbf{Accuracy} & \textbf{F1} & \textbf{ROC-AUC} & \textbf{MCC} \\
\hline
Random Forest & 0.78 & 0.72 & 0.81 & 0.58 \\
\hline
Gradient Boosting & 0.80 & 0.74 & 0.83 & 0.61 \\
\hline
LightGBM (Simple) & 0.80 & 0.73 & 0.84 & 0.62 \\
\hline
\textbf{LightGBM (Tuned)} & \textbf{0.82} & \textbf{0.76} & \textbf{0.86} & \textbf{0.65} \\
\hline
\end{tabular}
}
\end{center}

\vspace{0.3cm}
\textbf{Key Observations:}

\begin{itemize}
    \item Tuning improves \textbf{all 4 metrics}
    \item GB $>$ RF consistently (sequential $>$ parallel)
    \item LightGBM competitive or better than GB
    \item No single model dominates all metrics
    \item Gap between best and worst: $\sim$4-5\% ROC-AUC (clinically significant!)
\end{itemize}

\end{frame}

%=============================================================
% SLIDE 19: Common Mistakes
%=============================================================
\begin{frame}[fragile]
\frametitle{Hyperparameter Tuning: Mistakes to Avoid}

\textbf{Mistake 1: Tuning on test set}
\begin{itemize}
    \item \textcolor{red}{\ding{56}} Leads to overfitting to test data
    \item \textcolor{green}{\checkmark} Use cross-validation on training set
\end{itemize}

\textbf{Mistake 2: Too few trials}
\begin{itemize}
    \item \textcolor{red}{\ding{56}} 5-10 trials is random luck
    \item \textcolor{green}{\checkmark} Use 50+ trials
\end{itemize}

\textbf{Mistake 3: Ignoring min\_data\_in\_leaf}
\begin{itemize}
    \item \textcolor{red}{\ding{56}} LightGBM overfits on small data (n=299)
    \item \textcolor{green}{\checkmark} Set min\_data\_in\_leaf $\geq$ 20
\end{itemize}

\textbf{Mistake 4: Optimizing accuracy only}
\begin{itemize}
    \item \textcolor{red}{\ding{56}} Ignores class imbalance (10\% positive)
    \item \textcolor{green}{\checkmark} Optimize ROC-AUC, report all 4 metrics
\end{itemize}

\textbf{Mistake 5: Different metrics for train/test}
\begin{itemize}
    \item \textcolor{red}{\ding{56}} Confuses signal from noise
    \item \textcolor{green}{\checkmark} Use same metric everywhere
\end{itemize}

\end{frame}

%=============================================================
% SLIDE 20: Exercise 1
%=============================================================
\begin{frame}[fragile]
\frametitle{Exercise 1: Learning Rate Impact}

\textbf{Task:} Compare two learning rates

\begin{itemize}
    \item \textbf{Model A:} learning\_rate=0.1 (current)
    \item \textbf{Model B:} learning\_rate=0.05 (slower)
    \item Keep: n\_estimators=100, max\_depth=3
\end{itemize}

\vspace{0.3cm}
\textbf{What to Compare:}

\begin{itemize}
    \item Which model is more accurate?
    \item Which has better F1-score?
    \item Which has higher ROC-AUC?
\end{itemize}

\vspace{0.3cm}
\textbf{Expected Insight:}

\begin{itemize}
    \item lr=0.1: Faster convergence, might overshoot optimum
    \item lr=0.05: Slower, more stable, often better on test set
    \item Small data (n=299) usually favors slower learning
\end{itemize}

\vspace{0.2cm}
\textbf{Your Task:} Run both models and report all 4 metrics for each.

\end{frame}

%=============================================================
% SLIDE 21: Exercise 2
%=============================================================
\begin{frame}[fragile]
\frametitle{Exercise 2: Tuning Speed vs Quality}

\textbf{Task:} Compare tuning budget impact

\begin{itemize}
    \item \textbf{Version 1:} n\_trials=20
    \begin{itemize}
        \item Time: $\sim$2 minutes
        \item Expected ROC-AUC: 0.83
    \end{itemize}

    \vspace{0.2cm}
    \item \textbf{Version 2:} n\_trials=50 (current)
    \begin{itemize}
        \item Time: $\sim$5 minutes
        \item Expected ROC-AUC: 0.86
    \end{itemize}

    \vspace{0.2cm}
    \item \textbf{Version 3:} n\_trials=100 (if brave)
    \begin{itemize}
        \item Time: $\sim$10 minutes
        \item Expected ROC-AUC: 0.865
    \end{itemize}
\end{itemize}

\vspace{0.3cm}
\textbf{Your Task:}

\begin{enumerate}
    \item Run Optuna with n\_trials=20
    \item Compare best hyperparameters vs n\_trials=50
    \item Assess: Is the extra 3 minutes worth 0.03 ROC-AUC gain?
\end{enumerate}

\vspace{0.2cm}
\textbf{Expected Learning:} Diminishing returns in optimization

\end{frame}

%=============================================================
% SLIDE 22: Exercise 3
%=============================================================
\begin{frame}[fragile]
\frametitle{Exercise 3: Choose Your Own Adventure}

\textbf{Select ONE path:}

\vspace{0.2cm}
\textbf{Option A: Aggressive Tuning (Wider Search)}
\begin{itemize}
    \item Expand ranges: learning\_rate 0.001-0.3, num\_leaves 31-100
    \item Explore extreme values
    \item Best if: You have computational time and like exploration
\end{itemize}

\vspace{0.2cm}
\textbf{Option B: Manual Theory-Based Tuning}
\begin{itemize}
    \item Set hyperparameters from learned principles
    \item learning\_rate=0.05, num\_leaves=15 (conservative)
    \item Best if: You want to apply lessons learned
\end{itemize}

\vspace{0.2cm}
\textbf{Option C: Ultra-Conservative Model}
\begin{itemize}
    \item Optimize for stability over accuracy
    \item Very strict regularization (min\_data\_in\_leaf=30)
    \item Best if: You're concerned about overfitting
\end{itemize}

\vspace{0.3cm}
\textbf{Your Task:}

\begin{enumerate}
    \item Choose A, B, or C
    \item Run the code
    \item Write one paragraph observation comparing to LightGBM Tuned
\end{enumerate}

\end{frame}

%=============================================================
% SLIDE 23: Bias-Variance Tradeoff
%=============================================================
\begin{frame}[fragile]
\frametitle{Building Intuition: Bias-Variance Tradeoff}

\vspace{-0.3cm}
\begin{center}
Model Complexity $\rightarrow$
\end{center}

\textbf{High Variance (Overfitting)} \hfill \textbf{Low Variance (Underfitting)}

\vspace{0.1cm}
\begin{itemize}
    \item Shallow trees \hfill Deep trees
    \item Slow learning rate \hfill Fast learning rate
    \item Small n\_estimators \hfill Large n\_estimators
    \item High regularization \hfill No regularization
\end{itemize}

\vspace{0.3cm}
\textbf{Where models fall on this spectrum:}

\begin{itemize}
    \item \textbf{Random Forest:} Deep trees, many parallel \\
    $\rightarrow$ Low bias, \textcolor{red}{high variance}

    \item \textbf{Gradient Boosting:} Shallow trees, sequential \\
    $\rightarrow$ Medium bias, medium variance

    \item \textbf{LightGBM Tuned:} Leaf-wise optimization \\
    $\rightarrow$ \textcolor{green}{Low bias, low variance} \checkmark

    \item \textbf{LightGBM Conservative:} Strict rules \\
    $\rightarrow$ \textcolor{red}{High bias}, very low variance
\end{itemize}

\end{frame}

%=============================================================
% SLIDE 24: LightGBM Decision Making
%=============================================================
\begin{frame}[fragile]
\frametitle{How LightGBM Actually Makes Decisions}

\textbf{Decision at each expansion step:}

{\small
\begin{lstlisting}
When expanding a tree:

1. Current state: 10 leaves (potential splits)

2. Find the leaf with highest loss reduction:
   - Leaf A: Splitting reduces loss by 0.05
   - Leaf B: Splitting reduces loss by 0.12  ← Choose this
   - Leaf C: Splitting reduces loss by 0.03

3. Smart focus on high-error regions:
   - LightGBM: Splits Leaf B first (0.12 > others)
   - Traditional GB: Splits ALL leaves at depth d
                     equally (wastes effort)

4. Result:
   - LightGBM: Fewer leaves, faster training
   - Traditional: More leaves, slower, same/worse accuracy
\end{lstlisting}
}

\vspace{0.2cm}
\textbf{Smaller trees $\rightarrow$ Faster training + Better generalization}

\end{frame}

%=============================================================
% SLIDE 25: Real-World Context
%=============================================================
\begin{frame}[fragile]
\frametitle{Real-World Application: Building a Survival Predictor}

\textbf{Our Dataset:}

\begin{itemize}
    \item \textbf{299 heart failure patients}
    \item \textbf{11 clinical measurements} (age, ejection fraction, etc.)
    \item \textbf{Target:} Mortality (10\% death rate)
    \item \textbf{Imbalance:} Heavy class imbalance
\end{itemize}

\vspace{0.3cm}
\textbf{Algorithm Choice Rationale:}

\begin{center}
{\small
\begin{tabular}{|l|c|l|}
\hline
\textbf{When} & \textbf{Algorithm} & \textbf{Why} \\
\hline
Baseline & Random Forest & Easy to understand, fast, good benchmark \\
\hline
Better accuracy & Gradient Boosting & Sequential learning corrects RF errors \\
\hline
Production & LightGBM & Fast, tunable, handles constraints \\
\hline
\end{tabular}
}
\end{center}

\vspace{0.3cm}
\textbf{Real Constraints:}

\begin{itemize}
    \item \textbf{Small dataset} (n=299) $\rightarrow$ need strong regularization
    \item \textbf{Imbalanced classes} (10\% positive) $\rightarrow$ need multiple metrics
    \item \textbf{Medical domain} $\rightarrow$ need interpretable results
\end{itemize}

\end{frame}

%=============================================================
% SLIDE 26: Chicco & Jurman Connection
%=============================================================
\begin{frame}[fragile]
\frametitle{Connecting to Chicco \& Jurman (2020)}

\textbf{The Paper's Main Finding:}

\begin{quote}
For imbalanced binary classification, \textbf{MCC} is the most robust metric \\
that provides a single quality measure valid for both classes.
\end{quote}

\vspace{0.3cm}
\textbf{How Week 6 Applies This:}

\begin{itemize}
    \item[\checkmark] Use heart failure data (imbalanced: 10\% positive)

    \item[\checkmark] Report all 4 metrics (not accuracy alone)

    \item[\checkmark] Compare multiple models (RF, GB, LightGBM)

    \item[\checkmark] Use cross-validation (honest, unbiased estimates)

    \item[\checkmark] Follow ensemble best practices
\end{itemize}

\vspace{0.3cm}
\textbf{The Key Insight:}

\begin{itemize}
    \item Accuracy alone misleads on imbalanced data
    \item MCC catches weak models that accuracy misses
    \item Ensemble methods beat single models consistently
    \item Proper tuning gains: +0.04-0.05 ROC-AUC is clinically meaningful
\end{itemize}

\end{frame}

%=============================================================
% SLIDE 27: Key Takeaways
%=============================================================
\begin{frame}[fragile]
\frametitle{Key Takeaways}

\textbf{1. Algorithm Progression}
\begin{itemize}
    \item Random Forest: Parallel, simple, good baseline
    \item Gradient Boosting: Sequential, better, more complex
    \item LightGBM: Optimized, practical, fastest
\end{itemize}

\textbf{2. Why Sequential $>$ Parallel}
\begin{itemize}
    \item GB learns from residuals $\rightarrow$ reduces bias + variance
    \item RF only reduces variance via averaging
\end{itemize}

\textbf{3. Metrics Matter}
\begin{itemize}
    \item Accuracy alone is misleading
    \item MCC and ROC-AUC catch hidden problems
    \item Report all four, compare all
\end{itemize}

\textbf{4. Hyperparameter Tuning}
\begin{itemize}
    \item Optuna: Smart search saves time
    \item Cross-validation: Essential for honest estimates
    \item Regularization: Prevent overfitting on small data
\end{itemize}

\textbf{5. Production Readiness}
\begin{itemize}
    \item LightGBM is 10-20× faster
    \item Tuning gives consistent +0.03-0.05 gains
    \item Small data needs conservative settings
\end{itemize}

\end{frame}

%=============================================================
% SLIDE 28: What's Next
%=============================================================
\begin{frame}[fragile]
\frametitle{What's Next (Weeks 7+)}

\textbf{You now know:}

\begin{itemize}
    \item[\checkmark] Compare different algorithms systematically
    \item[\checkmark] Evaluate with multiple metrics
    \item[\checkmark] Tune hyperparameters automatically
    \item[\checkmark] Build ensemble models
\end{itemize}

\vspace{0.3cm}
\textbf{Coming next:}

\begin{itemize}
    \item Advanced feature engineering
    \item Handling class imbalance (SMOTE, class\_weight)
    \item Model interpretation (SHAP values, feature importance)
    \item Production deployment (model saving, serving)
    \item Cross-validation strategies for small data
\end{itemize}

\vspace{0.3cm}
\textbf{The big picture:}

Systematic, data-driven approach beats guessing every time.

\end{frame}

%=============================================================
% SLIDE 29: Resources
%=============================================================
\begin{frame}[fragile]
\frametitle{Resources for Deeper Learning}

\textbf{Gradient Boosting Intuition}
\begin{itemize}
    \item StatQuest: Gradient Boost (YouTube)
    \item StatQuest: XGBoost (YouTube playlist)
\end{itemize}

\textbf{LightGBM Mastery}
\begin{itemize}
    \item LightGBM Official Documentation
    \item Parameter Tuning Guide
\end{itemize}

\textbf{Hyperparameter Optimization}
\begin{itemize}
    \item Optuna Official Documentation
    \item Hyperparameter Optimization Strategies Tutorial
\end{itemize}

\textbf{Evaluation Best Practices}
\begin{itemize}
    \item Chicco \& Jurman (2020): MCC Advantages
    \item Imbalanced-Learn Library
\end{itemize}

\textbf{Live Examples}
\begin{itemize}
    \item Kaggle competitions using LightGBM
    \item GitHub: LightGBM competition winners
\end{itemize}

\end{frame}

%=============================================================
% SLIDE 30: Q&A
%=============================================================
\begin{frame}[fragile]
\frametitle{Q\&A / Open Discussion}

\textbf{Q: Why not just use the simplest model?}
\begin{itemize}
    \item A: Tuning gives +0.03-0.05 ROC-AUC on this dataset. In medical contexts, this can save lives.
\end{itemize}

\textbf{Q: Is LightGBM always better?}
\begin{itemize}
    \item A: Not always. On very small data (<100 samples) or with outliers, simpler models work better.
\end{itemize}

\textbf{Q: How many trials do I need?}
\begin{itemize}
    \item A: 50 is good for this problem. More trials rarely help beyond n=100 (diminishing returns).
\end{itemize}

\textbf{Q: What if my tuned model is worse on test data?}
\begin{itemize}
    \item A: Use CV-based tuning on training set. Report test-set results separately.
\end{itemize}

\textbf{Q: Should I report accuracy or ROC-AUC?}
\begin{itemize}
    \item A: Both. Report all four metrics. ROC-AUC more reliable for imbalanced data.
\end{itemize}

\end{frame}

%=============================================================
% SLIDE 31: The Big Picture
%=============================================================
\begin{frame}[fragile]
\frametitle{The Big Picture: From Week 1 to Week 6}

\vspace{-0.2cm}
\textbf{Progression:}

\begin{center}
{\small
\begin{tabular}{|l|l|}
\hline
Week 1-2 & Data loading + cleaning \\
\hline
Week 3-4 & Exploratory analysis + feature engineering \\
\hline
Week 5 & Baseline modeling (Random Forest) \\
\hline
\textbf{Week 6} & \textbf{\textcolor{red}{Systematic improvement (GB + LightGBM)}} \\
\hline
Week 7+ & Advanced techniques + production \\
\hline
\end{tabular}
}
\end{center}

\vspace{0.3cm}
\textbf{What this progression teaches:}

\begin{enumerate}
    \item \textbf{Start simple} (RF: 5 hyperparameters)
    \item \textbf{Improve systematically} (GB: understand residuals)
    \item \textbf{Optimize carefully} (LightGBM: 7 hyperparameters + Optuna)
    \item \textbf{Evaluate honestly} (4 metrics + cross-validation)
    \item \textbf{Scale confidently} (ready for production)
\end{enumerate}

\vspace{0.3cm}
\begin{center}
\large
\textit{Good machine learning is iterative, not magical.}

\vspace{0.2cm}
Simple baselines $\rightarrow$ Systematic improvements $\rightarrow$ Production systems
\end{center}

\end{frame}

%=============================================================
% FINAL SLIDE
%=============================================================
\begin{frame}[fragile]
\frametitle{Thank You}

\begin{center}
\Large
\textbf{Week 6: From Random Forest to LightGBM}

\vspace{0.5cm}
\normalsize
Questions? Comments? Observations?

\vspace{0.7cm}
\textit{Remember: The best model is the one you understand \\ and can explain to others.}

\end{center}

\end{frame}

\end{document}
